% Listati della pipeline sperimentale. I sorgenti eseguibili restano l'autorità;
% gli estratti mostrano il contratto versione 2 senza incorporare risultati.

\begin{onehalfspace}

Questa sezione raccoglie il materiale implementativo trasferito dal
Capitolo~\ref{chap:implementazione}: contratto di esecuzione, organizzazione dei moduli,
dettagli dei codici QEC ed estratti dei sorgenti. Gli output numerici non sono codificati
nei listati: ogni campagna li salva in un artefatto JSON schema~2 corredato di manifest.

\subsection{Contratto implementativo e riproducibilità}
\label{app:contratto_codice}

L'ambiente Python è isolato in un virtual environment su \ac{WSL}; versioni di
interprete, Qiskit, Aer, NumPy, SciPy e scikit-learn vengono registrate nel manifest di
ogni artefatto. La campagna canonica usa Qiskit~2.5 e simulazione \ac{CPU}. Il
Listato~\ref{lst:installazione_dellambiente_python_e} documenta la creazione
dell'ambiente, mentre il Listato~\ref{lst:errore_gpu_su_wsl} conserva l'esito negativo
del tentativo GPU: l'accelerazione non è quindi parte del contratto riproducibile.

\subsubsection{Netlist compilata e manifest}
\label{app:netlist_manifest}

La compilazione è centralizzata e condivisa da training, baseline e sweep. Il contratto
fissa \texttt{optimization\_level=2}, seed del transpiler e base nativa
\texttt{rz/sx/x/cx}. Con Qiskit~2.5, $N=15$, $a=7$ e otto qubit di conteggio, la
netlist ha profondità 412 e contiene 224 porte \texttt{cx}. Per $N=21$, $a=2$, la
variante Beauregard raggiunge invece profondità 23\,081 e 21\,036 \texttt{cx}; a
$\lambda=10^{-3}$ il proxy registrato vale $2.6987\times10^{-9}$. Quest'ultimo numero
documenta il costo circuitale, non la probabilità di successo, ed è la ragione pratica
per cui N21/N35 restano nelle validazioni ideali.

Per ogni compilazione il manifest conserva dimensione, profondità, conteggio completo
delle operazioni, base, livello e seed di transpiler, versione delle librerie e SHA-256
della sequenza compilata. Modelli e risultati dichiarano inoltre schema, revisione,
preset, seed e definizione dell'etichetta; l'inferenza rifiuta un classificatore con
manifest incompatibile. In tal modo un artefatto non può essere associato per errore a
una revisione differente del circuito.

\subsubsection{Seed, tie-break e confronto appaiato}

Per le campagne N15, una sola simulazione fornisce l'istogramma comune a M1, TOP-4 e M2.
Il seed è separato fra repliche e iterazioni secondo
\begin{equation}
\texttt{seed\_simulator}=\texttt{rep}\cdot1\,000\,000+
\texttt{iteration}\cdot10\,000,
\end{equation}
una distanza superiore al numero di shot che evita la quasi sovrapposizione degli stream
Aer osservata con semi consecutivi. Nei pareggi, una priorità SHA-256 calcolata da seed e
bitstring evita dipendenze dall'ordine del dizionario o dal valore numerico del candidato.

La prima iterazione di successo è registrata separatamente per ciascuna strategia; un
fallimento entro \texttt{max\_iter} è codificato, ai soli fini inferenziali, con
\texttt{max\_iter+1}. Il Wilcoxon appaiato usa
\texttt{zero\_method='pratt'}: test unilaterale per M1$>$TOP-4 e M1$>$M2, bilaterale
per TOP-4--M2. Il JSON schema~2 conserva input completi, vettori appaiati, seed, test,
versioni software, hash e manifest; figure e tabelle leggono questi artefatti senza
trascrizione manuale.

\subsection{Organizzazione completa del codice sperimentale}
\label{app:organizzazione_codice}

I sorgenti eseguibili sono suddivisi per campagna, così che codice, README, log e
artefatti restino associati allo stesso contratto. La Tabella~\ref{tab:organizzazione_qec}
esplicita la mappa usata nel testo; le sigle M identificano le milestone sperimentali e
non livelli di affidabilità del risultato.

\begin{table}[H]
\centering
\small
\begin{tabularx}{\textwidth}{XXXX}
\toprule
\textbf{Cartella} & \textbf{Oggetto} & \textbf{Strumento} & \textbf{Costo dominante} \\
\midrule
\path{campagne_classiche_M1-M4} & Shor, TOP-K e classificatore & Qiskit Aer + scikit-learn & traiettorie rumorose \\
\path{M5_repetition_code} & codice a ripetizione & Qiskit Aer & trascurabile \\
\path{M6_steane_code} & Steane $[[7,1,3]]$ & Qiskit Aer & Monte Carlo Pauli \\
\path{M7_surface_code} & surface code & Stim + PyMatching & decodifica \\
\path{M8_shor_logico} & proxy per gate e TOP-4 M13 & Qiskit Aer & matrix product state \\
\path{M10_neural_decoder} & decodifica appresa e ibrida & Stim + PyMatching + scikit-learn & addestramento \\
\path{M11_layout} & layout su snapshot eterogenea & Qiskit Aer & compilazioni replicate \\
\path{M12_coerenti} & controllo su errori coerenti & Qiskit Aer & simulazione full-state \\
\path{M15_qft_topologia} & QFT approssimata & Qiskit Aer & sweep di compilazione \\
\path{extra_rumore_coerente} & Pauli contro sovrarotazione & Qiskit Aer & controllo full-state 1Q \\
\bottomrule
\end{tabularx}
\caption[Organizzazione del codice sperimentale]{Organizzazione completa delle campagne.
La toolchain cambia con la classe del circuito: Aer per circuiti generali di taglia
trattabile, Stim e PyMatching per circuiti stabilizer su larga scala.}
\label{tab:organizzazione_qec}
\end{table}

\subsection{Dettagli implementativi dei componenti QEC}
\label{app:dettagli_qec}

\subsubsection{M5: codice a ripetizione}

Gli stabilizzatori sono misurati indirettamente copiandone la parità su qubit ancillari,
senza misurare i dati e collassare l'informazione logica. Porte identità collocate nel
punto di iniezione costituiscono l'aggancio a cui Aer associa il canale senza modificare
lo stato ideale. Le varianti in base $Z$ e $X$ condividono lo stesso schema, con i cambi
di base necessari alla rilevazione, rispettivamente, di bit-flip e phase-flip
(Listato~\ref{lst:qec_repetition}).

\subsubsection{M6: codice di Steane}

Lo stato $\ket{0_L}$ è preparato ponendo tre qubit seme in sovrapposizione e
propagandone il supporto con CNOT. Gli stabilizzatori CSS di tipo $X$ e $Z$ alimentano
una lookup table di correzione a peso minimo. La stima di $p_L$ usa un modello di Pauli;
l'estrazione di sindrome è ideale e non costituisce un protocollo fault-tolerant
completo (Listato~\ref{lst:qec_steane}).

\subsubsection{M7: surface code}

Stim genera il circuito di memoria e campiona i detection event; PyMatching costruisce
dal detector error model il decoder \ac{MWPM}. Nell'esperimento sul crosstalk, il canale
fra dati adiacenti è inserito ricorsivamente anche nei blocchi \texttt{REPEAT}; i qubit
di misura sono identificati dal primo \texttt{MR}, non dalla misura finale che comprende
anche i dati. I Listati~\ref{lst:crosstalk} e~\ref{lst:stim_pymatching} mostrano i due
passaggi.

\subsubsection{M8/M13 e M10: proxy di Shor e decoder appreso}

Nel proxy a livelli, $p_L$ è la probabilità totale del Pauli non-identità applicato dopo
ogni gate incluso nel modello. Prima dell'esecuzione Aer viene convertito in
$\lambda_1=4p_L/3$ per i gate 1Q e $\lambda_2=16p_L/15$ per quelli 2Q; le \texttt{rz}
restano virtuali e i canali agiscono su \texttt{sx/x} e \texttt{cx}. Questo parametro
per gate non coincide con il logical error rate per round o memoria di M6/M7 senza una
compilazione fault-tolerant e un mapping esplicito. In M10, decoder appreso e
\ac{MWPM} ricevono esattamente gli stessi detection event, così il confronto riguarda la
regola di decodifica e non i campioni osservati.

\end{onehalfspace}

\subsection{Estratti dei sorgenti}

\begin{lstlisting}[caption={Installazione dell'ambiente Python e delle librerie necessarie}, label={lst:installazione_dellambiente_python_e}]
python3 -m venv ~/quantum-env
source ~/quantum-env/bin/activate
python -m pip install -r requirements.txt
python -c "import qiskit, qiskit_aer; print(qiskit.__version__)"
\end{lstlisting}

\begin{lstlisting}[caption={Errore GPU osservato su WSL durante i test}, label={lst:errore_gpu_su_wsl}]
RuntimeError: Simulation device "GPU" is not supported on this system
\end{lstlisting}

\begin{lstlisting}[caption={Implementazione della Quantum Fourier Transform inversa}, label={lst:implementazione_della_quantum_fourier}]
def inverse_qft(n_qubits):
    qc = QuantumCircuit(n_qubits, name='QFT_dagger')
    for q in range(n_qubits // 2):
        qc.swap(q, n_qubits - q - 1)
    for j in range(n_qubits):
        for m in range(j):
            qc.cp(-np.pi / 2 ** (j - m), m, j)
        qc.h(j)
    return qc
\end{lstlisting}

\begin{lstlisting}[caption={Moltiplicazione modulare N15 corretta: l'operazione $\times a$ viene ripetuta \texttt{power} volte}, label={lst:modular_exponentiation_per_n15}]
def c_amod15(a, power):
    if a not in [2, 7, 8, 11, 13]:
        raise ValueError("base non supportata per N=15")
    U = QuantumCircuit(4)
    for _ in range(power):
        if a in [2, 13]:
            U.swap(0, 1); U.swap(1, 2); U.swap(2, 3)
        if a in [7, 8]:
            U.swap(2, 3); U.swap(1, 2); U.swap(0, 1)
        if a == 11:
            U.swap(1, 3); U.swap(0, 2)
        if a in [7, 11, 13]:
            U.x(range(4))
    return U.control(annotated=False)
\end{lstlisting}

\begin{lstlisting}[caption={Circuito N15 e post-processing classico}, label={lst:circuito_completo_di_shor}]
def shor_circuit(N, a, n_count):
    if N != 15:
        return shor_circuit_beauregard(N, a, n_count)
    n_work = 4
    qc = QuantumCircuit(n_count + n_work, n_count)
    qc.h(range(n_count))
    qc.x(n_count + 3)   # |1> nella convenzione del registro textbook
    for j in range(n_count):
        power = 2 ** j
        if power % 4 != 0:       # U^4 = I per a=7 mod 15
            qc.append(c_amod15(a, power),
                      [j] + list(range(n_count, n_count + n_work)))
    qc.barrier()
    qc.append(inverse_qft(n_count), range(n_count))
    qc.measure(range(n_count), range(n_count))
    return qc

def extract_factors(measured_value, n_count, N, a):
    if measured_value == 0:
        return None, None
    phase = measured_value / 2 ** n_count
    r = Fraction(phase).limit_denominator(N).denominator
    if r % 2:
        return None, None
    for candidate in (gcd(a ** (r // 2) - 1, N),
                      gcd(a ** (r // 2) + 1, N)):
        if 1 < candidate < N:
            return candidate, N // candidate
    return None, None
\end{lstlisting}

\begin{lstlisting}[caption={Noise model uniforme illustrativo: RZ virtuale, durate 1Q/2Q separate}, label={lst:costruzione_del_noise_model}]
BASIS_GATES = ('rz', 'sx', 'x', 'cx')
GATE_TIME_1Q_NS = 50.0
GATE_TIME_2Q_NS = 300.0

def compose_errors(errors):
    out = None
    for error in errors:
        if error is not None:
            out = error if out is None else out.compose(error)
    return out

def build_noise_model(eps_1q, eps_2q, t1_ns, t2_ns, p_ro=0.02):
    # eps_1q ed eps_2q sono i lambda di depolarizing_error.
    thermal_1q = thermal_relaxation_error(
        t1_ns, t2_ns, GATE_TIME_1Q_NS)
    thermal_one_2q = thermal_relaxation_error(
        t1_ns, t2_ns, GATE_TIME_2Q_NS)
    thermal_2q = thermal_one_2q.tensor(thermal_one_2q)
    error_1q = compose_errors([
        depolarizing_error(eps_1q, 1) if eps_1q else None,
        thermal_1q,
    ])
    error_2q = compose_errors([
        depolarizing_error(eps_2q, 2) if eps_2q else None,
        thermal_2q,
    ])
    nm = NoiseModel()
    nm.add_all_qubit_quantum_error(error_1q, ['sx', 'x'])
    nm.add_all_qubit_quantum_error(error_2q, ['cx'])
    if p_ro:
        nm.add_all_qubit_readout_error(
            ReadoutError([[1-p_ro, p_ro], [p_ro, 1-p_ro]]))
    return nm
\end{lstlisting}

\begin{lstlisting}[caption={Compilazione deterministica, tie-break SHA-256 e baseline TOP-1}, label={lst:pipeline_del_metodo_1}]
def compile_shor_circuit(N, a, n_count):
    return transpile(
        shor_circuit(N, a, n_count),
        basis_gates=['rz', 'sx', 'x', 'cx'],
        optimization_level=2,
        seed_transpiler=20260819,
    )

def rank_measurements(counts, tie_seed):
    def priority(bitstring):
        payload = f'{int(tie_seed)}:{bitstring}'.encode('ascii')
        return hashlib.sha256(payload).digest()
    return sorted(counts.items(),
                  key=lambda item: (-item[1], priority(item[0])))

def method1_from_counts(counts, tie_seed, n_count, N, a):
    bitstring = rank_measurements(counts, tie_seed)[0][0]
    return extract_factors(int(bitstring, 2), n_count, N, a)
\end{lstlisting}

\begin{lstlisting}[caption={Dataset condiviso per l'audit delle etichette TOP-1 e TOP-16}, label={lst:generazione_del_dataset_con}]
LABEL_TOP_KS = (1, 16)
labels = {top_k: [] for top_k in LABEL_TOP_KS}

for sample_index in range(n_samples):
    sample_seed = seed * 1_000_000 + sample_index * 10_000
    counts = simulator.run(
        compiled, noise_model=sample_noise, shots=shots,
        seed_simulator=sample_seed,
    ).result().get_counts()
    ranked = rank_measurements(counts, tie_seed=sample_seed)
    for top_k in LABEL_TOP_KS:
        useful = any(
            extract_factors(int(bits, 2), n_count, N, a)[0] is not None
            for bits, _ in ranked[:top_k]
        )
        labels[top_k].append(int(useful))
    X.append(normalized_histogram(counts, n_count, shots))

# Ogni modello salva schema_version=2.0, label_top_k e manifest.
\end{lstlisting}

\begin{lstlisting}[caption={Decisione M2 sullo stesso istogramma usato da M1 e TOP-4}, label={lst:loop_di_inferenza_del}]
def evaluate_same_histogram(counts, simulation_seed, classifier,
                            n_count, N, a, shots):
    ranked = rank_measurements(counts, simulation_seed)
    valid = [
        extract_factors(int(bits, 2), n_count, N, a)[0] is not None
        for bits, _ in ranked[:4]
    ]
    m1_success = valid[0]
    top4_success = any(valid)
    feature = normalized_histogram(counts, n_count, shots)
    m2_success = bool(classifier.predict([feature])[0]) and top4_success
    return m1_success, top4_success, m2_success
\end{lstlisting}

\begin{lstlisting}[caption={Orchestrazione appaiata delle sole campagne rumorose N15}, label={lst:schema_di_orchestrazione_degli}]
PRESETS = {
    'reference': dict(N=15, a=7, n_count=8,
                      eps_1q=1e-3, eps_2q=1e-2,
                      t1_ns=100_000, t2_ns=80_000, p_ro=0.02),
    'stress':    dict(N=15, a=7, n_count=8,
                      eps_1q=5e-3, eps_2q=5e-2,
                      t1_ns=50_000, t2_ns=30_000, p_ro=0.05),
}
K_REPETITIONS, SHOTS, MAX_ITER = 30, 1024, 50

for rep in range(K_REPETITIONS):
    first = {'m1': None, 'top4': None, 'm2': None}
    for iteration in range(1, MAX_ITER + 1):
        simulation_seed = rep * 1_000_000 + iteration * 10_000
        counts = simulator.run(compiled, shots=SHOTS,
             seed_simulator=simulation_seed).result().get_counts()
        flags = evaluate_same_histogram(counts, simulation_seed, ...)
        update_first_success(first, flags, iteration)
    encoded = {name: value if value is not None else MAX_ITER + 1
               for name, value in first.items()}

# Inferenza: wilcoxon(paired_a, paired_b,
#                     zero_method='pratt', alternative=...)
# Output: schema_version='2.0', input espliciti e manifest/hash.
\end{lstlisting}

\begin{lstlisting}[caption={Codifica ed estrazione di sindrome del codice a ripetizione.}, label={lst:qec_repetition}]
if logical_value == 1:
    qc.x(0)
qc.cx(0, 1); qc.cx(0, 2)
if basis == 'X':
    qc.h(DATA)
for q in DATA:
    qc.id(q)                       # punto di iniezione del rumore
if basis == 'Z':
    qc.cx(0, 3); qc.cx(1, 3)
    qc.cx(1, 4); qc.cx(2, 4)
else:
    qc.h(3); qc.cx(3, 0); qc.cx(3, 1); qc.h(3)
    qc.h(4); qc.cx(4, 1); qc.cx(4, 2); qc.h(4)
qc.measure(ANC[0], 0); qc.measure(ANC[1], 1)
\end{lstlisting}

\begin{lstlisting}[caption={Preparazione di $\ket{0_L}$ e misura degli stabilizzatori di Steane.}, label={lst:qec_steane}]
def encode_zero_L(qc):
    for seed in SEED:
        qc.h(seed)
    for seed, support in SEED.items():
        for data in support:
            if data != seed:
                qc.cx(seed, data)

def measure_syndrome(qc):
    for k, support in enumerate(GZ):
        for data in support:
            qc.cx(data, ANC_Z[k])
    for k, support in enumerate(GX):
        qc.h(ANC_X[k])
        for data in support:
            qc.cx(ANC_X[k], data)
        qc.h(ANC_X[k])
\end{lstlisting}

\begin{lstlisting}[caption={Iniezione ricorsiva del crosstalk nel circuito del surface code.}, label={lst:crosstalk}]
def inject(block, pairs, p_ct):
    out = stim.Circuit()
    first = True
    for inst in block:
        if isinstance(inst, stim.CircuitRepeatBlock):
            out += inject(inst.body_copy(), pairs, p_ct) * inst.repeat_count
            continue
        out.append(inst)
        if inst.name == 'DEPOLARIZE1':
            if first:
                out.append('DEPOLARIZE2', pairs, p_ct)
            first = not first
    return out
\end{lstlisting}

\begin{lstlisting}[caption={Stima di $p_L$ con Stim e PyMatching.}, label={lst:stim_pymatching}]
circuit = stim.Circuit.generated(
    'surface_code:rotated_memory_x', distance=d, rounds=rounds,
    after_clifford_depolarization=p,
    after_reset_flip_probability=p,
    before_measure_flip_probability=p,
)
sampler = circuit.compile_detector_sampler()
events, observables = sampler.sample(shots, separate_observables=True)
model = circuit.detector_error_model(decompose_errors=True)
matching = pymatching.Matching.from_detector_error_model(model)
predictions = matching.decode_batch(events)
p_L = (predictions != observables).any(axis=1).mean()
\end{lstlisting}
